\documentclass[pdflatex,sn-mathphys-num]{sn-jnl}

\usepackage{graphicx}%
\usepackage{multirow}%
\usepackage{amsmath,amssymb,amsfonts}%
\usepackage{amsthm}%
\usepackage{mathrsfs}%
\usepackage[title]{appendix}%
\usepackage{xcolor}%
\usepackage{textcomp}%
\usepackage{manyfoot}%
\usepackage{booktabs}%
\usepackage{algorithm}%
\usepackage{algorithmicx}%
\usepackage{algpseudocode}%
\usepackage{listings}%
\usepackage{lmodern}%
\usepackage{siunitx}
\usepackage{microtype}
\usepackage{derivative}
\usepackage[version=4]{mhchem}
\usepackage{cleveref}
\usepackage[spanish]{babel}

\renewcommand{\arraystretch}{1.2}
\sisetup{group-digits=true,
  group-separator={\,},
  separate-uncertainty
}

\theoremstyle{thmstyleone}%
\newtheorem{theorem}{Teorema}%
\newtheorem{proposition}[theorem]{Proposición}%

\theoremstyle{thmstyletwo}%
\newtheorem{example}{Ejemplo}%
\newtheorem{remark}{Observación}%

\theoremstyle{thmstylethree}%
\newtheorem{definition}{Definición}%

\DeclareSIUnit\angstrom{\text{\r{A}}}

\raggedbottom
%%\unnumbered% uncomment this for unnumbered level heads

\begin{document}

\title[Física de Plasmones]{Física de Plasmones de Superficie y Biosensores:
Resumen de la Charla del Prof. César Herreño}

\author{\fnm{Julian L.}
\sur{Avila-Martinez}}\email{jlavilam@udistrital.edu.co}

\affil{\orgdiv{Programa Académico de Física}, \orgname{Universidad Distrital
Francisco José de Caldas}, \orgaddress{\city{Bogotá D.C.},
\country{Colombia}}}

\abstract{
Se resume la formulación física y las aplicaciones de la
resonancia de plasmones de superficie presentadas por César Herreño.
Se distinguen los mecanismos de confinamiento radiativo y la dispersión modal
para excitaciones localizadas y propagantes. Se analizan las
implicaciones del acoplamiento difractivo en redes periódicas y se establecen
las aplicaciones clínicas y medioambientales de la transducción plasmónica.
}

\keywords{Plasmónica, relaciones de dispersión, biosensores, anomalías de
Rayleigh, termodinámica de nanopartículas}

\maketitle

\section{Electrodinámica de Plasmones y Dispersión}
\label{sec:plasmones}

La excitación colectiva cuantizada del gas de electrones exhibe dos regímenes
fundamentales determinados por las condiciones de frontera. El primer
régimen, la resonancia de plasmón de superficie localizado o LSPR, emerge
cuando la radiación electromagnética polariza una nanopartícula e induce un
dipolo oscilante. Esta oscilación, confinada espacialmente por la
morfología superficial, amplifica el campo eléctrico en varios órdenes de
magnitud dentro de volúmenes nanométricos denominados puntos calientes.
La disipación térmica de la energía electromagnética absorbida ocurre mediante
radiación de frenado o bremsstrahlung originada por el dipolo.

El segundo régimen comprende los polaritones de plasmón de superficie, o SPP,
definidos como excitaciones propagantes a lo largo de una interfaz
dieléctrico-metal. En el espacio libre surge una restricción
cinemática fundamental debido a la linealidad de la relación de dispersión del
fotón incidente frente a la no linealidad de la curva de dispersión del
plasmón. Esta asimetría impide la conservación del momento bajo iluminación
directa. La configuración de Kretschmann resuelve este desacople al
modular el vector de onda del fotón mediante una interfaz de dos materiales
dieléctricos distintos, satisfaciendo así el ángulo crítico necesario para
excitar el modo SPP.

\section{Resonancia de Redes y Sensado Optoelectrónico}
\label{sec:redes}

La maximización de la sensibilidad analítica en biosensores exige una
estructuración espacial precisa de las nanopartículas. Simulaciones
computacionales sobre un dominio de \qty{300}{\nano\meter} revelan que la
respuesta electromagnética de una configuración en heptámero depende
críticamente de la distancia de separación o gap entre los constituyentes
individuales. La organización de estas entidades en bases de red
periódicas como discos, esferas o varillas permite un acoplamiento
difractivo de los modos LSPR mediado por anomalías de Rayleigh, induciendo la
resonancia de red de superficie, SLR. Este fenómeno colectivo
reduce sustancialmente el ancho de línea espectral y eleva la sensibilidad en
varios órdenes de magnitud respecto a los sistemas LSPR aislados.

\section{Aplicaciones Biológicas y Perspectivas Clínicas}
\label{sec:aplicaciones}

La espectroscopía SPR permite la cuantificación directa de las afinidades de
enlace molecular. Mediciones experimentales de la interacción entre
el rotavirus y proteínas receptoras específicas demuestran que la
funcionalización proteica con polietilenglicol, PEG, incrementa notablemente
la sección eficaz de unión viral.

Las aplicaciones tecnológicas de estas configuraciones incluyen la síntesis de
calcogenuros y óxidos de hierro para remediación ambiental, enfocándose en la
adsorción de metales pesados dentro de ecosistemas fluviales.
Paralelamente, se investigan rutas de síntesis verde para nanopartículas de
plata mediante el uso de biomasa de \textit{Cannabis sativa}. En
oncología, el mecanismo intrínseco de disipación térmica del régimen LSPR
fundamenta el diseño de protocolos clínicos que integran radiación ionizante
con hipertermia localizada.

\end{document}
